\documentclass[10pt,a4paper]{amsart}
\usepackage{amsmath,amssymb,amsthm,graphicx}
% For commutative diagrams we use Xy-pic
\usepackage[all]{xy}

% We use 2cell for 2-commutative diagrams.
\xyoption{2cell}
% Theorem environments.
%
\theoremstyle{plain}
\newtheorem{theorem}[subsection]{Theorem}
\newtheorem{proposition}[subsection]{Proposition}
\newtheorem{lemma}[subsection]{Lemma}

\theoremstyle{definition}
\newtheorem{definition}[subsection]{Definition}
\newtheorem{example}[subsection]{Example}
\newtheorem{exercise}[subsection]{Exercise}
\newtheorem{situation}[subsection]{Situation}

\theoremstyle{remark}
\newtheorem{remark}[subsection]{Remark}
\newtheorem{remarks}[subsection]{Remarks}

\numberwithin{equation}{subsection}

% Macros
%
\def\lim{\mathop{\mathrm{lim}}\nolimits}
\def\colim{\mathop{\mathrm{colim}}\nolimits}
\def\Spec{\mathop{\mathrm{Spec}}}
\def\Hom{\mathop{\mathrm{Hom}}\nolimits}
\def\Ext{\mathop{\mathrm{Ext}}\nolimits}
\def\SheafHom{\mathop{\mathcal{H}\!\mathit{om}}\nolimits}
\def\SheafExt{\mathop{\mathcal{E}\!\mathit{xt}}\nolimits}
\def\Sch{\mathit{Sch}}
\def\Mor{\mathop{\mathrm{Mor}}\nolimits}
\def\Ob{\mathop{\mathrm{Ob}}\nolimits}
\def\Sh{\mathop{\mathit{Sh}}\nolimits}
\def\NL{\mathop{N\!L}\nolimits}
\def\CH{\mathop{\mathrm{CH}}\nolimits}
\def\proetale{{pro\text{-}\acute{e}tale}}
\def\etale{{\acute{e}tale}}
\def\QCoh{\mathit{QCoh}}
\def\Ker{\mathop{\mathrm{Ker}}}
\def\Im{\mathop{\mathrm{Im}}}
\def\Coker{\mathop{\mathrm{Coker}}}
\def\Coim{\mathop{\mathrm{Coim}}}

% Boxtimes
%
\DeclareMathSymbol{\boxtimes}{\mathbin}{AMSa}{"02}

%
% Macros for moduli stacks/spaces
%
\def\QCohstack{\mathcal{QC}\!\mathit{oh}}
\def\Cohstack{\mathcal{C}\!\mathit{oh}}
\def\Spacesstack{\mathcal{S}\!\mathit{paces}}
\def\Quotfunctor{\mathrm{Quot}}
\def\Hilbfunctor{\mathrm{Hilb}}
\def\Curvesstack{\mathcal{C}\!\mathit{urves}}
\def\Polarizedstack{\mathcal{P}\!\mathit{olarized}}
\def\Complexesstack{\mathcal{C}\!\mathit{omplexes}}
% \Pic is the operator that assigns to X its picard group, usage \Pic(X)
% \Picardstack_{X/B} denotes the Picard stack of X over B
% \Picardfunctor_{X/B} denotes the Picard functor of X over B
\def\Pic{\mathop{\mathrm{Pic}}\nolimits}
\def\Picardstack{\mathcal{P}\!\mathit{ic}}
\def\Picardfunctor{\mathrm{Pic}}
\def\Deformationcategory{\mathcal{D}\!\mathit{ef}}
\usepackage{cleveref}
\crefname{originalreference}{original reference}{original references}
\Crefname{originalreference}{Original reference}{Original references}
\makeatletter
\newlabel{lemma-directed-colimit-cohomology}{{03Q6}{source}{}{sourcepointer-1}{}}
\newlabel{lemma-directed-colimit-cohomology@cref}{{[originalreference][0][]03Q6}{[1][0][]source}{}{}{}}
\makeatother
\crefname{originalreference}{original reference}{original references}
\Crefname{originalreference}{Original reference}{Original references}
\makeatletter
\newlabel{lemma-vanishing-etale-cohomology-strictly-henselian}{{03QO}{source}{}{sourcepointer-2}{}}
\newlabel{lemma-vanishing-etale-cohomology-strictly-henselian@cref}{{[originalreference][0][]03QO}{[1][0][]source}{}{}{}}
\makeatother
\title{If all local rings are strictly henselian then etale and Zariski cohomology agree}
\author{The Stacks Project Authors}
\date{stacks-project:etale-cohomology.tex:local-snapshot-20261011; GNU FDL 1.2 or later}
\begin{document}
\maketitle
\noindent\textit{Editorial note.} All delivered mathematical text is the author's own, copied from the pinned Stacks Project chapter file etale-cohomology.tex; no statement or proof step was rewritten, reordered or removed.\par\medskip
\noindent\textit{Editorial note.} The counted proof is the author's complete argument: the proof is written out in full: the description of the higher direct images, the identification of the stalk through the colimit lemma, and the vanishing over strictly henselian local rings leading to the Leray degeneration.\par\medskip
\noindent\textit{Editorial note.} The label lemma-local-rings-strictly-henselian names the counted statement, and the proof environment immediately following it in the source is delivered verbatim as its proof.\par\medskip
\noindent\textit{Editorial note.} The counted statement is carried by the author's own lemma environment in the chapter source.\par\medskip
\noindent\textit{Editorial note.} The statement and proof are copied from the pinned chapter file; the two cited inputs are kept as external pointers with their canonical Stacks tags.\par\medskip
\noindent\textit{Allowed standard prerequisites (editorial).}\begin{itemize}
\item Higher direct images of the comparison morphism as sheafified cohomology. Same chapter, the description of R\textasciicircum{}p epsilon\_*F as the sheaf associated to U |-> H\textasciicircum{}p\_etale(U, F). Supplies the computation of the stalks of the higher direct images.
\item Cohomology of a filtered colimit of schemes. Same chapter, the lemma computing the stalk of the sheafified cohomology as the cohomology of Spec of the local ring with pulled-back coefficients. Supplies the identification of the stalk as H\textasciicircum{}p\_etale(Spec(O\_\{X,x\}), G\_x).
\item Vanishing of etale cohomology over a strictly henselian local ring. Same chapter, the lemma stating the vanishing of the higher etale cohomology of a strictly henselian local ring. Supplies the vanishing that makes the higher direct images zero.
\end{itemize}\medskip
\ifdefined\hypertarget\hypertarget{sourcepointer-1}{}\fi
\noindent\textit{Source reference.} Original reference 03Q6: sources/stacks-etale-cohomology/pinned-original-source.tex, source lines 7464--7464. Prerequisite pointer invoked at proof line(s) 14220; it names a standard fact of the same chapter that the counted argument uses, and every step of the counted proof built on it is retained verbatim in the delivered body.\par\medskip
\ifdefined\hypertarget\hypertarget{sourcepointer-2}{}\fi
\noindent\textit{Source reference.} Original reference 03QO: sources/stacks-etale-cohomology/pinned-original-source.tex, source lines 8207--8207. Prerequisite pointer invoked at proof line(s) 14223; it names a standard fact of the same chapter that the counted argument uses, and every step of the counted proof built on it is retained verbatim in the delivered body.\par\medskip
\begin{lemma}
\label{lemma-local-rings-strictly-henselian}
Let $S$ be a scheme all of whose local rings are strictly henselian.
Then for any abelian sheaf $\mathcal{F}$ on $S_\etale$ we have
$H^i(S_\etale, \mathcal{F}) = H^i(S_{Zar}, \mathcal{F})$.
\end{lemma}

\begin{proof}
Let $\epsilon : S_\etale \to S_{Zar}$ be the morphism of sites given
by the inclusion functor. The Zariski sheaf $R^p\epsilon_*\mathcal{F}$
is the sheaf associated to the presheaf $U \mapsto H^p_\etale(U, \mathcal{F})$.
Thus the stalk at $x \in X$ is
$\colim H^p_\etale(U, \mathcal{F}) =
H^p_\etale(\Spec(\mathcal{O}_{X, x}), \mathcal{G}_x)$
where $\mathcal{G}_x$ denotes the pullback of $\mathcal{F}$
to $\Spec(\mathcal{O}_{X, x})$, see
Lemma \ref{lemma-directed-colimit-cohomology}.
Thus the higher direct images of $R^p\epsilon_*\mathcal{F}$ are
zero by
Lemma \ref{lemma-vanishing-etale-cohomology-strictly-henselian}
and we conclude by the Leray spectral sequence.
\end{proof}

\end{document}
